\section{Torres en la máscara residual}\label{sec:torres}
\label{sec:rook}
La reducción a un permanente auxiliar, usada para el término constante~\cite{tc}, permite comparar
dos tableros de llegada. En el primero se conservan las casillas disponibles y se levantan las
restricciones diagonales entre las piezas que faltan. Sus completaciones son los emparejamientos
perfectos de la máscara $M$, y su suma ponderada es $\operatorname{per}M$. En el segundo fijamos
dos torres en casillas $e,f$ de una diagonal. Las piezas restantes ocupan la máscara $M_{ef}$,
que pierde dos filas y dos columnas. El cociente $w_ew_f\operatorname{per}M_{ef}/\operatorname{per}M$
da la probabilidad de ocupar el par. Al aproximar ambos permanentes, ¿cuánto del error del primer
tablero se cancela en el segundo? En la Figura~\ref{fig:mascara-ilustrada} se puede leer el
numerador a mano para $e=(2,4)$ y $f=(3,3)$: sólo queda la asignación $(4,1)$. Así, uno de los
dos emparejamientos de la máscara contiene ese par.

\paragraph{Clase y dos operaciones.}
\[
  w_e\in\{0\}\cup[w_0,w_1],\quad P=\operatorname{diag}(u)M\operatorname{diag}(v),
  \quad P_e\le M_0/k,\quad \|P-J_k\|_{\rm op}\le1-\gamma .
\]
Aquí $0<w_0\le w_1$, $M_0\ge1$, $0<\gamma\le1$, $M$ tiene soporte total, $P$ es su escalado
bistocástico y $J_k$ tiene entradas $1/k$. Ésta es la clase $\mathcal K$. En esta sección $k$ es
el tamaño entero de la máscara, distinto de la función $k(x)$ de la integral. Ocupar $e$ retira
su fila y su columna y deja el menor $M_e$. Prohibir $e$ deja las dos fibras y pone $w_e=0$ en
$M\setminus e$. Las dos operaciones perturban de manera distinta las sumas que se han de reescalar.
El punto estacionario $p$ de torres satisface $p_e(1-p_e)=w_e\exp(z_i+y_j)$
en la rama pequeña y tiene sumas de línea uno. Con $p$ se calcula el predictor logarítmico $L_T$
y con $P$, la corrección a $L_T$ dada por la fórmula siguiente.
Las cotas siguientes valen para $k$ mayor que un umbral de clase, también en los menores utilizados.
El corte físico se fija antes de enviar $n$ a infinito.

\paragraph{La aproximación que se transporta [D].}
\[
   g_T(M)=\log\operatorname{per}M-L_T(M)
       =\frac{\Phi(kP)}k+R(M),\qquad |R(M)|\le Ck^{-2}.
\]
Esta estimación de valor es (T). Al escribir $\Phi=\mathfrak F-1/12$ y evaluar $\mathfrak F$ en
el perfil de torres de la trayectoria, se obtiene el $F(x)$ de~\eqref{eq:main}. Aplicar la cota
para $R$ por separado en los dos tableros sólo daría $O(k^{-2})$ para su diferencia. La estimación
$(T^+)$ gana un orden bajo $M\mapsto M_e$, con incremento $O(k^{-3})$, y dos bajo
$M\mapsto M\setminus e$, con incremento $O(k^{-4})$. Las constantes y el umbral de estas tres
cotas dependen sólo de $M_0,\gamma$. El apéndice~\ref{app:rook} obtiene las ganancias siguiendo
por separado el resto del permanente escalado y el de la normalización del predictor.

\paragraph{De dos torres a un par de reinas.}
En una máscara binaria, la probabilidad exacta de ocupar $E=\{e,f\}$, con filas y columnas
distintas, es $\pi^{\rm ex}=\operatorname{per}M_{ef}/\operatorname{per}M$.
Para una única línea de dos casillas, el término correspondiente en el logaritmo del conteo es
$\log(1-\pi^{\rm ex})$. La aproximación local $\pi_1$ se define en el
apéndice~\ref{app:pair}. Si se revela una tercera casilla $c$ en otras filas y columnas,
la variación del error $\pi^{\rm ex}-\pi_1$ cuesta a lo sumo
$C\varpi(E)/k^5$, donde $\varpi(E)=k^2p_ep_f$. Este enunciado (R1) se aplica al menor de
torres que conserva el par. En la geometría física una jugada también puede vedar casillas
adicionales. Al tomar la media de llegada, el cambio del error de torres aporta $P_T$. El efecto
de una sola línea contribuye a $J$, y las dos diagonales de un cuadrado completo producen, además,
el término $-S$ del apéndice~\ref{app:boundary}.
